\section{Finite monoids: verified obstructions and counterexamples}\label{sec:A}

\subsection{A1: the two proofs of the finite-dimensional obstruction}

\begin{lemma}[A1.1: the group at the minimal ideal]\label{A:minimal-group}
Let $M$ be a finite monoid. Its universal group $U(M)$ is finite, the canonical
homomorphism $\eta:M\to U(M)$ is surjective, and $\pi_1(BM)=U(M)$.
If $e$ is an idempotent in the minimal two-sided ideal of $M$, then $H=eMe$ is a
finite group and $\eta|_H:H\to U(M)$ is surjective. There is a map
$j_e:BH\to BM$, defined through the Karoubi envelope, which induces this
surjection on fundamental groups.
\end{lemma}
\begin{proof}
\emph{Step 1, the presentation and finiteness.}
The one-vertex nerve has an edge for each nonidentity element and a triangular
relation $[a][b]=[ab]$ for every pair. This is exactly the presentation of the
universal group. For each $a\in M$ choose $r<s$ with $a^r=a^s$. In the group,
$\eta(a)^{s-r}=1$, so $\eta(a)^{-1}=\eta(a^{s-r-1})$. Thus the image of $M$ is
already closed under inverses and is all of $U(M)$, of order at most $|M|$.

\emph{Step 2, the corner.}
A finite monoid has a nonempty minimal ideal $K$: the intersection of its
nonempty two-sided ideals is nonempty, since the product of any two such ideals
is contained in their intersection. A suitable positive power of any element
of $K$ is an idempotent $e$. For $x\in eMe$, minimality gives $MxM=K$ and hence
$e=axb=(eae)x(ebe)$. In a finite monoid, an element with a one-sided inverse is
invertible: its corresponding multiplication map is injective or surjective,
hence bijective. If a product is the identity, each factor is consequently a
unit. Applied in $eMe$, this makes every $x$ invertible. Finally $\eta(e)=1$,
and $\eta(eme)=\eta(m)$ proves surjectivity of the corner map.

\emph{Step 3, typing the map.}
The category $\operatorname{Kar}(M)$ has idempotents as objects and
$\operatorname{Hom}(f,g)=gMf$, with multiplication as composition. Inclusion of
the object $1$ is a functor $i:M\to\operatorname{Kar}(M)$. For an object $e$,
an object of $i/e$ is $m\in eM$, and a morphism $m\to n$ is $a\in M$ with
$na=m$. The arrows labelled $m$ give a natural transformation from the identity
of $i/e$ to the constant functor at $e$: indeed $em=m$ and $na=m$. Thus $B(i/e)$
is contractible. Quillen's Theorem A \cite[Section~1]{quillen1973} implies
$Bi$ is a homotopy equivalence.
The unital functor $H\to\operatorname{Kar}(M)$ selecting $e$, followed by a
homotopy inverse of $Bi$, defines $j_e$. The two arrows labelled $e$ between
$1$ and $e$ transport a loop $h$ to the edge $h$ of $BM$. This proves the stated
map on $\pi_1$, without treating the nonunital inclusion $eMe\to M$ as a functor.
\end{proof}

\begin{theorem}[A1.2: finite-dimensional obstruction]\label{A:finite-dimension}
If $M$ is finite and $BM$ has the homotopy type of a finite-dimensional CW
complex, then $U(M)=1$.
\end{theorem}
\begin{proof}
\emph{Step 4, Miller's input and its hypotheses.}
Miller's Theorem B, stated on page 40 and proved as Theorem 10.1 on pages
84--85 of \emph{The Sullivan conjecture on maps from classifying spaces}
(1984) \cite{miller1984}, states that a map from the classifying space of a torsion group to a
finite-dimensional CW complex induces the zero homomorphism on $\pi_1$.
Here the source group is the finite group $H$ of Lemma~\ref{A:minimal-group};
the target is the given finite-dimensional model $X$ of $BM$; the map is
$BH\xrightarrow{j_e}BM\simeq X$. The 1985 correction \cite{miller1985} was read in full. Its two
repairs concern the unstable-coalgebra spectral sequence and Fact 3.9; they
change neither the statement nor the independent K-theory proof in Section 10.

\emph{Step 5, the conclusion.}
The homomorphism $H\to\pi_1(X)$ is simultaneously surjective by
Lemma~\ref{A:minimal-group} and zero by the preceding paragraph. Therefore
$\pi_1(X)=U(M)=1$. These five steps confirm the five steps of
\texttt{G6\_VERDICT.md}; the comma-category contraction above supplies the
comparison explicitly.
\end{proof}

\begin{proposition}[A1.3: the K-theory proof independently of Miller]\label{A:atiyah}
If $G$ is a torsion group, $X$ is a finite-dimensional CW complex, and
$f:BG\to X$ is a map, then $f_*=0$.
\end{proposition}
\begin{proof}
Suppose some $g\in G$ has nontrivial image. A primary component of $g$ gives a
cyclic $p$-group $C\leq G$ whose image $N$ is a nontrivial cyclic $p$-group.
The map $BC\to X$ lifts to the covering $X_N$ corresponding to $N\leq\pi_1X$.
The cover still has dimension $d=\dim X$. Choose a character $\psi:N\to S^1$
of order $p$, and let $L$ be its flat line bundle on $X_N$. The reduced class
$v=[L]-1$ has skeletal filtration at least one, so $v^{d+1}=0$.

The pullback to $BC$ is the reduced character $\chi-1$, with $\chi$ of order
$p$. Its $(d+1)$st power therefore restricts to zero on every finite skeleton
of $BC$. Atiyah's Theorem (7.2) identifies the inverse limit of these K-groups
with the augmentation-adic completion of $R(C)$; Proposition (6.11) says that
$R(C)$ injects into that completion for a $p$-group. Both statements, and the
more precise cyclic Proposition (8.1), were read in the attached 1961 paper.
Thus $(\chi-1)^{d+1}=0$ in $R(C)$. Evaluation at an element on which $\chi$ is a
primitive $p$th root $\zeta$ gives $(\zeta-1)^{d+1}=0$ in $\mathbb C$, a
contradiction. This also checks the cyclic K-theory calculation behind
Miller's Section 10, without relying on the corrected portions of his paper
\cite{atiyah1961}.
\end{proof}

\subsection{A2: cohomological bounds, filtrations and finite homotopy type}

\begin{lemma}[A2.1: skeletal filtration with bounded ordinary cohomology]\label{A:filtration}
For a CW complex $Y$, put
$F^qK^0(Y)=\ker(K^0(Y)\to K^0(Y^{q-1}))$.
The filtration is multiplicative and natural for cellular maps.
If $H^q(Y;\mathbb Z)=0$ for $q>n$, then
$F^{n+1}K^0(Y)=\bigcap_q F^qK^0(Y)$.
The same statement holds for the spectrum $KU_{(p)}$ and coefficients
$\mathbb Z_{(p)}$.
\end{lemma}
\begin{proof}
For multiplication, lift classes of filtrations $a,b$ to the relative groups
for $(Y,Y^{a-1})$ and $(Y,Y^{b-1})$. Their external product vanishes on
$Y^{a-1}\times Y\cup Y\times Y^{b-1}$. A cellular diagonal sends the
$(a+b-1)$-skeleton into this union, proving $F^aF^b\subset F^{a+b}$.
Naturality follows from the defining kernels after cellular approximation.

Here is the important stabilization step, avoiding any assumption of strong
convergence of an infinite Atiyah--Hirzebruch spectral sequence. If
$x\in F^qK^0(Y)$, its restriction to $Y^q$ is the image of a relative class
$c\in K^0(Y^q,Y^{q-1})$. For odd $q$ this relative group is zero. For even $q$
it is the group of cellular $q$-cochains with integer coefficients. Since $x$
extends over $Y^{q+1}$, the connecting homomorphism for the triple of skeleta
annihilates $c$; under the cellular identification this says $\delta c=0$.
For $q>n$, write $c=\delta b$. The class $b$ belongs to
$K^{-1}(Y^{q-1},Y^{q-2})$. Its image in $K^{-1}(Y^{q-1})$, followed by the
connecting map, is $c$, so exactness of the pair $(Y^q,Y^{q-1})$ makes the
image of $c$ in $K^0(Y^q)$ zero. Hence $x\in F^{q+1}$. This proves
$F^q=F^{q+1}$ for all $q>n$, and therefore stabilization to the intersection.
For $KU_{(p)}$ the identical argument has coefficient groups
$\mathbb Z_{(p)}$ in even degree and zero in odd degree.
\end{proof}

\begin{theorem}[A2.2: representation obstruction]\label{A:representation}
Let $X$ be connected with $H^q(X;\mathbb Z)=0$ for $q>n$, and let $G$ be finite.
For every map $f:BG\to X$, the subgroup $f_*(G)$ lies in
\[
 R_{\mathbb C}(\pi_1X)=
 \bigcap_{d\geq1}\ \bigcap_{\rho:\pi_1X\to\operatorname{GL}_d(\mathbb C)}
 \ker\rho.
\]
In particular $f_*=0$ if $\pi_1X$ is finite, residually finite, or a complex
linear group. If $\pi_1X$ is finite and $H^q(X;\mathbb Z_{(p)})=0$ for $q>n$,
then $|f_*(G)|$ is prime to $p$.
\end{theorem}
\begin{proof}
\emph{The flat bundle and the integral argument.}
Let $\rho:\pi_1X\to\operatorname{GL}_d(\mathbb C)$ be any finite-dimensional
representation, and let $V$ be the flat bundle induced from the universal
cover. The rank-zero class $v=[V]-d$ has filtration at least one. By
Lemma~\ref{A:filtration}, $v^{n+1}$ restricts to zero on every skeleton of $X$.
Make $f$ cellular. On every finite skeleton of $BG$, the pullback of this class
is zero. It is the image of the virtual representation
$(\rho f_*-d)^{n+1}$ under Atiyah's completion map. Theorem (7.2) and
Proposition (6.10) of Atiyah say that this virtual character vanishes on every
prime-power-order element of $G$. Thus
$\operatorname{tr}(\rho f_*(g))=d$ for each such $g$. A finite-order complex
matrix is diagonalizable with root-of-unity eigenvalues. Their sum equals
$d$ only when each eigenvalue is one, so $\rho f_*(g)=1$. Every element of a
finite group is a product of its prime-power primary components. This proves
the representation-radical assertion. Finite groups have faithful regular
representations; a finite quotient separates each nontrivial element of a
residually finite group; and a complex linear group has a faithful defining
representation. These observations check all three stated corollaries.

\emph{The $p$-local assertion.}
Use the regular representation of $\pi=\pi_1X$, of degree $d=|\pi|$.
If $p$ divides $|f_*(G)|$, choose $u\in f_*(G)$ of order $p$. A lift in $G$
has a $p$-primary component $c$ still mapping to $u$. Restrict to the cyclic
$p$-group $C=\langle c\rangle$. The localized version of the preceding
filtration argument kills $(\rho f_*|_C-d)^{n+1}$ in $KU_{(p)}^0$ of every
finite odd-dimensional lens-space skeleton of $BC$. Its reduced integral
$K^0$ is a finite $p$-group: the finite Atiyah--Hirzebruch filtration has
finite $p$-groups in every positive even filtration and zero otherwise.
Consequently localization at $p$ is injective on these reduced K-groups.
The class is zero integrally on every skeleton. Atiyah (7.2), (6.11) now makes
it zero in $R(C)$, whereas evaluation at $c$ is $(-d)^{n+1}\ne0$. This is the
required contradiction. All uses of finite skeleta, rather than an unproved
identification of localized infinite K-groups, are essential here.
\end{proof}

\begin{theorem}[A2.3: the precise bounded-homology theorem]\label{A:bounded}
For a finite monoid $M$ the following are equivalent: $BM$ has finite CW
homotopy type; it has finite-dimensional CW homotopy type; and
$H_k(M;\mathbb Z)=0$ for all sufficiently large $k$. Each implies $U(M)=1$.
If $p\mid |U(M)|$, then $H_k(M;\mathbb F_p)\ne0$ for infinitely many $k$.
A finite projective resolution of the trivial left or right $\mathbb ZM$-module
therefore forces finite simply connected homotopy type.
\end{theorem}
\begin{proof}
The first two implications are immediate. The normalized nerve has finitely
many cells in each dimension, so all integral homology groups are finitely
generated. A high-degree integral homology bound gives a high-degree integral
cohomology bound by the universal coefficient exact sequence. Apply
Theorem~\ref{A:representation} to $j_e:BH\to BM$: its target has finite
fundamental group by Lemma~\ref{A:minimal-group}, so its surjection on
fundamental groups is zero. Thus $BM$ is simply connected.

Hatcher's Proposition 4C.1 \cite[pp.~429--430]{hatcher2002}, read in the attached text,
constructs a CW model for any simply connected CW complex whose homology is
a direct sum of cyclic groups, with one generator cell per cyclic summand
and one next-dimensional relation cell for every finite cyclic summand.
Here the groups are finitely generated and only finitely many are nonzero,
so this model is finite (of dimension at most one above the highest nonzero
homology degree). The proposition does not assert that the attaching maps
are determined by homology; only the numbers of these cells are used.

If instead mod-$p$ homology vanishes in all large degrees, the universal
coefficient sequence and finite generation imply that integral homology
localized at $p$ vanishes there. Applying the localized conclusion of
Theorem~\ref{A:representation} to $j_e$ contradicts $p\mid |U(M)|$.
Finally a finite projective resolution makes the relevant Tor groups zero
above its length. These arguments confirm the individual steps of
\texttt{OWN\_RESULTS.md}, Theorems 1.1--1.2 and 2.1--2.4.
The assertion 2.5 additionally depends on the positive realization theorem
checked separately in Part B.
\end{proof}

\subsection{A3: two proposed strengthenings fail, and a six-element cup-product obstruction}

\begin{theorem}[A3.1: simply connected does not imply nilpotent positive cohomology]\label{A:six-nonnilpotent}
There is a six-element finite regular monoid $M$ with $U(M)=1$ and a class
$u\in H^2(M;\mathbb F_2)$ such that $u^k\ne0$ for every $k\geq1$.
This is class 190 of the six candidate tables in
\texttt{higher\_homology\_order6.out}.
\end{theorem}
\begin{proof}
Let $M=\{1,g\}\amalg\{x_{ij}:i,j\in\mathbb F_2\}$. Declare $g^2=1$, and
\[
 x_{ij}x_{kl}=x_{il},\qquad
 gx_{ij}=x_{i+1,j},\qquad x_{ij}g=x_{i,j+1}.
\]
The rectangular-band rule is associative; left and right flips commute and
respect the appropriate outer coordinate; two flips are the identity.
These facts check all placements of group and band factors in an associative
triple. In the order $(1,g,x_{00},x_{10},x_{01},x_{11})$, the full table is
\[
\begin{array}{c|rrrrrr}
 &1&g&x_{00}&x_{10}&x_{01}&x_{11}\\\hline
1&1&g&x_{00}&x_{10}&x_{01}&x_{11}\\
g&g&1&x_{10}&x_{00}&x_{11}&x_{01}\\
x_{00}&x_{00}&x_{01}&x_{00}&x_{00}&x_{01}&x_{01}\\
x_{10}&x_{10}&x_{11}&x_{10}&x_{10}&x_{11}&x_{11}\\
x_{01}&x_{01}&x_{00}&x_{00}&x_{00}&x_{01}&x_{01}\\
x_{11}&x_{11}&x_{10}&x_{10}&x_{10}&x_{11}&x_{11}
\end{array}
\]
All band elements are idempotent, and $g$ is invertible, so $M$ is regular.
Every idempotent maps to one in a group; the equation $gx_{00}=x_{10}$ then
kills $g$. Hence $U(M)=1$.

Define normalized one-cochains $F,H:M\to\mathbb F_2$ by
\[
 F(1)=H(1)=0,\quad F(g)=H(g)=1,\quad
 F(x_{ij})=1+j,\quad H(x_{ij})=1+i.
\]
Write $\epsilon$ for the indicator of the band ideal. Direct substitution in
the multiplication rules gives
\[
 (\delta F)(a,b)=F(a)\epsilon(b),\qquad
 (\delta H)(b,c)=\epsilon(b)H(c).
\]
Therefore the two-cochain $u(a,b)=F(a)H(b)$ satisfies
$\delta u=(\delta F)\smile H-F\smile\delta H=0$.
The inclusion of the unit group $\{1,g\}=C_2$ is unital. On that subgroup
$F=H=t$, the normalized degree-one generator. Thus restriction sends
$[u]$ to $t^2\in H^*(C_2;\mathbb F_2)=\mathbb F_2[t]$ and sends $[u]^k$ to
$t^{2k}\ne0$. One can see the last fact directly: the normalized bar
cochain complex of $C_2$ over $\mathbb F_2$ has one basis cochain in every
degree and zero differential, and all these cup powers evaluate to one
on the word $(g,\ldots,g)$.

The independent program \texttt{code/a\_counterexamples.py} verifies all
$6^3=216$ associativity equations and all 216 cocycle equations. Its output
also checks the unit, the explicit universal-group-killing certificate,
and the restriction. The infinite nonvanishing assertion is proved by
restriction, not inferred from a finite range of computed powers.
\end{proof}

\begin{proposition}[A3.2: the requested six initial tests]\label{A:six-tests}
For the six labelled tables specified in the prompt, $H^2(-;\mathbb F_2)$ is
one-dimensional. The square of its nonzero class is zero for classes
$138,152,1285,1296,1317$, and nonzero for class $190$. Restriction to the
unit $C_2$ is zero in the first two cases and nonzero in class $190$.
\end{proposition}
\begin{proof}
The same independent program constructs normalized bar differentials over
$\mathbb F_2$ and performs bit-exact Gaussian elimination. It computes
$\ker(\delta:C^2\to C^3)/\operatorname{im}(\delta:C^1\to C^2)$, chooses a
nonzero representative, forms its Alexander--Whitney square, and tests
membership in $\operatorname{im}(\delta:C^3\to C^4)$. The complete tables,
representatives and Boolean outcomes are recorded in
\texttt{code/outputs/a\_counterexamples.json}; the text output states the
scope. This is a computation on six supplied tables, with all axioms
rechecked, and is not represented as an independent enumeration. Nilpotence
of all other classes is not inferred from these degree-two tests.
\end{proof}

\begin{proposition}[A3.3: bounded exponent does not constrain the universal group]\label{A:exponent}
The condition that $H_k(M;\mathbb Z)$ be finite of uniformly bounded exponent
for all sufficiently large $k$ does not imply $U(M)=1$. Indeed every finite
group occurs as $U(M)$ for a monoid satisfying this condition.
\end{proposition}
\begin{proof}
For $M=C_p=\langle g\mid g^p=1\rangle$, use the periodic free resolution over
$\mathbb ZC_p$ with alternating maps $g-1$ and
$N=1+g+\cdots+g^{p-1}$. If $a=\sum a_i g^i$, then
$(g-1)a=0$ exactly when all $a_i$ are equal, so its kernel is $\mathbb ZN$.
Also $Na=(\sum a_i)N$, whose kernel is the augmentation ideal, exactly the
image of multiplication by $g-1$. This proves exactness in all degrees.
Tensoring with the trivial module gives alternating maps $0$ and $p$, hence
\[
 H_{2j+1}(C_p;\mathbb Z)=\mathbb Z/p\quad(j\geq0),\qquad
 H_{2j}(C_p;\mathbb Z)=0\quad(j\geq1).
\]
The universal group is $C_p$. Tables and periodic matrices for
$p=2,3,5,7$ are checked in the program, while the argument proves every prime.

More generally take $M=Q$ for any finite group $Q$. The covering
$EQ\to BQ$ has $|Q|$ sheets. Transfer, defined by summing all lifts of each
singular simplex, composes with projection to multiplication by $|Q|$ on
homology. Since $EQ$ is contractible, $|Q|$ annihilates every positive
homology group. Each is finitely generated by the finite bar complex in each
degree, hence finite. Thus the hypothesis permits every possible finite
universal group; no group-theoretic restriction stronger than finiteness
can follow from it alone.
\end{proof}

\begin{theorem}[A3.4: no unrestricted-fundamental-group version of R4]\label{A:acyclic}
There is a connected integrally acyclic CW complex $X$ and a map
$BC_2\to X$ inducing an injection on fundamental groups. Thus R4 is false
without its hypothesis on the target fundamental group. In fact any given
group embeds in the fundamental group of such an acyclic $X$.
\end{theorem}
\begin{proof}
Here is an explicit construction and an acyclicity proof. For a group $B$,
form $D=B\times B$, then the two HNN extensions
\[
 E=\langle D,d\mid d^{-1}(b,1)d=(b,b)\ (b\in B)\rangle,
\quad
 m(B)=\langle E,s\mid s^{-1}(b,1)s=(1,b)\ (b\in B)\rangle.
\]
The associated maps identify isomorphic subgroups, so the HNN normal form
embeds $D$ in $E$ and $E$ in $m(B)$. Embed $B$ as $B\times1$. In $m(B)$ its
conjugate $B^s$ commutes with $B$, and $b^d=bb^s$. Set
$B_0=C_2$, $B_{r+1}=m(B_r)$ and $A=\bigcup_r B_r$ with these embeddings.
This specifies the countable group completely by generators and relations;
in particular the original involution stays nonidentity.

We give the homology argument because merely naming an acyclic envelope
would not check the needed property. Suppose $\phi:L\to B$ induces zero on
field homology in degrees $1,\ldots,n-1$. Define
$\kappa:B\times B\to m(B)$ by $\kappa(a,b)=ab^s$ and let
$\mu:B\to m(B)$ be inclusion. For $z\in H_n(L;k)$ the K\"unneth formula
and the vanishing assumption make the image of its diagonal equal to
$\phi_*(z)\otimes1+1\otimes\phi_*(z)$; all cross terms vanish.
The map $b\mapsto bb^s$ is conjugate to $\mu$, and both coordinate maps
followed by $\kappa$ are conjugate to $\mu$. Inner automorphisms act
identically on homology. Consequently
$(\mu\phi)_*(z)=2(\mu\phi)_*(z)$, and hence $(\mu\phi)_*(z)=0$ over every
field. Induction shows that a sufficiently long composite
$B_r\to B_{r+n}$ kills homology in all degrees through $n$.
Group homology commutes with directed unions, as every finite bar chain
uses only finitely many group elements. Thus $H_n(A;k)=0$ for all $n>0$
and every field $k$.

To pass to integral coefficients, $H_n(A;\mathbb Z)\otimes\mathbb Q=0$
makes it torsion. The universal coefficient sequence in degree $n+1$ and
$H_{n+1}(A;\mathbb F_p)=0$ show that it has no $p$-torsion for any prime.
It is therefore zero. Take $X=BA$. Universal coefficients also give
$H^q(X;\mathbb Z)=0$ for every $q>0$, whereas the initial inclusion
$C_2\hookrightarrow A$ gives the asserted map and injection.

This construction is the standard mitosis of Baumslag--Dyer--Heller \cite{baumslagdyerheller1980},
\emph{The topology of discrete groups}, equations (4.5)--(4.6),
Proposition 4.1, Theorem 4.2 and Lemma 5.4, all opened in the original paper.
The preceding argument verifies the hypotheses and derives the particular
counterexample. Replacing $C_2$ by any group gives the final assertion.
\end{proof}

\subsection{A5: what passes to finite categories}

\begin{theorem}[A5.1: local groups and endomorphism loops]\label{A:categories}
Let $C$ be a finite category. If $BC$ has finite-dimensional CW homotopy
type, every finite subgroup at every idempotent of every endomorphism monoid
maps trivially into $\pi_1(BC)$ through the Karoubi envelope. In fact every
endomorphism arrow of $C$ maps to the identity in the fundamental groupoid.
The same conclusion holds whenever $BC$ has bounded integral homology,
without any assumption on its component fundamental groups.
\end{theorem}
\begin{proof}
The inclusion $C\to\operatorname{Kar}(C)$ is a classifying-space equivalence:
for an idempotent $(c,e)$, the comma category has objects arrows
$f:x\to c$ with $ef=f$, and $f$ itself defines the natural transformation
from its identity to the constant object $(c,e)$. Every group
$G\leq e\operatorname{End}_C(c)e$ is a group of automorphisms of $(c,e)$ and
therefore supplies a well-typed map $BG\to B\operatorname{Kar}(C)\simeq BC$.
Miller's Theorem B applies in the finite-dimensional case.

For $N=\operatorname{End}_C(c)$, its minimal-ideal group surjects onto
$U(N)$ by Lemma~\ref{A:minimal-group}. The functor $N\to C$ selecting $c$
therefore kills every endomorphism loop once that group is killed.
For bounded integral homology, Theorem~\ref{A:representation} first gives
the conclusion for finite, residually finite, or complex linear fundamental
groups. The general case is established in
Theorem~\ref{A:bounded-categories} below, using the finite groupoid image
and an equivariant vector bundle; that proof uses none of the assertions
of the present theorem.
\end{proof}

\begin{proposition}[A5.2: the resulting fundamental groups have no further general restriction]\label{A:quotient}
Under either hypothesis of Theorem~\ref{A:categories}, there is a
finite acyclic category $D$ and a functor $C\to D$ inducing an equivalence
of fundamental groupoids. Nevertheless every finitely presented group is
$\pi_1(BC)$ for some finite category $C$ with finite-dimensional $BC$.
\end{proposition}
\begin{proof}
Quotient $C$ by the congruence imposing $a=1_c$ for every endomorphism
$a:c\to c$. This does not change objects or create arrows, and remains
finite. Its fundamental groupoid is the quotient by the normal closure of
the endomorphism loops, so Theorem~\ref{A:categories} makes the comparison
an equivalence. In this quotient, arrows in opposite directions between
two objects are mutual inverses. A skeleton is therefore acyclic: it has
only identity endomorphisms and no directed cycles. Choosing the skeleton
and the associated equivalence gives $D$ and the required functor.

Conversely, the finite two-dimensional presentation complex of any finite
group presentation has that fundamental group. Subdivide the finite
attaching loops and use simplicial mapping cylinders to replace the finite
CW complex by a finite simplicial complex of the same homotopy type.
Its nonempty-face poset is finite, and its nerve is the barycentric
subdivision. Taking this poset as $C$ proves the assertion. Thus neither
triviality, finiteness nor torsion-freeness of the fundamental group follows
for categories with several objects.
\end{proof}

\subsection{A3(i): nilpotence is detected on the local finite groups}

For an idempotent $e$ in a finite monoid, $G_e$ denotes the group of units
of $eMe$; the corner $eMe$ itself need not be a group away from the minimal
ideal. This distinction is necessary in the following theorem.

\begin{lemma}[A3.5: Sch\"utzenberger induction and apex filtration]\label{A:apex}
Let $M$ be a finite monoid, $k$ a field, and $e\in M$ idempotent. Set
\[
 I_e=\{m\in M:e\notin MmM\},\quad
 A_e=kM/kI_e,\quad L_e=Me\setminus I_e.
\]
Then $I_e$ is a two-sided ideal, $eA_ee=kG_e$, and $L_e$ is a free right
$G_e$-set. Consequently
\[
 \operatorname{Ind}_e W=kL_e\otimes_{kG_e}W
\]
is an exact functor from finite-dimensional $kG_e$-modules to
finite-dimensional $kM$-modules. Every simple $kM$-module $S$ has an
idempotent $e$, called an apex, for which $eS$ is a simple $kG_e$-module
and there is an exact sequence
\[
 0\longrightarrow K\longrightarrow\operatorname{Ind}_e(eS)
   \longrightarrow S\longrightarrow0.                 \tag{A.1}
\]
Every composition factor of $K$ has apex in a strictly larger regular
$\mathcal J$-class, where $e\leq_{\mathcal J}f$ means $e\in MfM$.
\end{lemma}
\begin{proof}
\emph{The corner and freeness.}
The complement condition makes $I_e$ an ideal. If $a\in eMe$ is outside
$I_e$, write $e=b a c=(ebe)a(ece)$. In the finite monoid $eMe$, a product
that is the identity has all factors invertible. Thus precisely the surviving
basis elements of $eA_ee$ are $G_e$.
For $l\in L_e$, choose $a,b\in M$ with $e=a l b$. Put $u=eal$ and $v=ebe$.
Then $u,v\in eMe$ and $uv=e$, so $u$ is invertible. The element
$t=u^{-1}ea$ satisfies $tl=e$. If $lh=lh'$ with $h,h'\in G_e$, multiplication
by $t$ gives $h=h'$. This proves freeness, without an appeal to an unstated
Green-stability theorem. The vector space $kL_e=A_ee$ has the quotient left
action: $m\cdot l=ml$ when $ml\in L_e$, and zero otherwise. Its freeness
on the right proves the exactness assertion.

\emph{Existence of an apex.}
Let $S$ be simple. Each two-sided ideal of $M$ either annihilates $S$ or has
$k$-linear span sending $S$ onto $S$. The product of two ideals of the latter
kind again has this property and is contained in their intersection.
There is therefore a unique smallest such ideal $J$, and $J^2=J$.
Write $I=\{m:mS=0\}$. The nonzero Rees quotient of $J$ by $J\cap I$
(adjoin a zero if that intersection is empty) has square equal to itself.
It has a nonzero idempotent. Here is the finite-semigroup argument: if every
nonzero element were nilpotent, a sufficiently long nonzero product would
have two equal nonzero prefixes, say $a=au$, implying $a=au^r=0$ for large
$r$. Thus the semigroup would be nilpotent, contrary to its square being
itself. An element that is not nilpotent has a nonzero idempotent power.
Lifting that idempotent gives $e\in J$ with $eS\ne0$.

For every $m$ that does not annihilate $S$, the ideal $MmM$ acts nontrivially,
so contains $J$ and hence $e$. Conversely, if $e\in MmM$, $m$ cannot annihilate
$S$. Thus $I=I_e$. The nonzero space $eS$ is simple over $eA_ee$: if
$0\ne W\subseteq eS$ is a submodule, then $A_eW=S$ by simplicity, and
$eA_eW=eA_eeW=W$, so $W=eS$.

\emph{The induction step.}
Multiplication gives the surjection in (A.1), since $A_eeS=S$.
Applying $e$ identifies its source and target with $eS$, so $eK=0$.
Moreover $I_eK=0$. If $T$ is a simple composition factor of $K$ with apex
$f$, then $fT\ne0$ and $I_eT=0$ imply $f\notin I_e$, that is
$e\leq_{\mathcal J}f$. Since $eT=0$, one has $e\in I_f$, so
$f\not\leq_{\mathcal J}e$. The inequality is strict, as asserted.
\end{proof}

\begin{lemma}[A3.6: the cohomology action on induced modules]\label{A:action-induction}
Let $x\in H^d(M;k)=\operatorname{Ext}_{kM}^d(k,k)$, and suppose its
restriction to $H^d(G_e;k)$ is zero. Then $x$ acts as zero on every
$\operatorname{Ind}_eW$ in the bounded derived category of $kM$-modules:
\[
 x\otimes\operatorname{id}_{\operatorname{Ind}_eW}:
 \operatorname{Ind}_eW\longrightarrow\operatorname{Ind}_eW[d]
 \quad\hbox{is zero}.
\]
The restriction meant here is the one induced by the local-group map
$BG_e\to B\operatorname{Kar}(M)\simeq BM$.
\end{lemma}
\begin{proof}
A $kM$-module $V$ extends canonically to the Karoubi envelope by assigning
$eV$ to $e$. This is an exact equivalence on module categories. Restriction
to $G_e$ is consequently the exact functor $V\mapsto eV$, and it carries
the constant module to the trivial $kG_e$-module. The standard bar resolution
identifies its map on Ext with the stated topological restriction.

Consider the two exact functors in the variable $V$,
\[
 T(V)=\operatorname{Ind}_e(W\otimes_k eV),\qquad
 U(V)=\operatorname{Ind}_eW\otimes_k V,
\]
with diagonal actions. There is a natural transformation
\[
 \Phi_V:T(V)\longrightarrow U(V),\qquad
 l\otimes(w\otimes v)\longmapsto(l\otimes w)\otimes lv.       \tag{A.2}
\]
It respects the $G_e$-balancing relation. It is $M$-linear: if $ml$ remains
in $L_e$, this is associativity; if $ml$ leaves $L_e$, both sides are zero
because the induced-module factor is zero. At $V=k$, (A.2) is the canonical
isomorphism. These facts do not claim that (A.2) is an isomorphism for
arbitrary $V$.

The exact functors and natural transformation pass to derived categories.
The morphism $T(x)$ is zero because $e(x)=\operatorname{res}_{G_e}(x)=0$.
Naturality and the isomorphism $\Phi_k$ then make $U(x)$ zero. Equivalently,
apply (A.2) to an extension representing $x$: it gives a morphism of
extensions which is an isomorphism at the two ends, so the Yoneda classes
agree, and the induced restricted extension has zero class.
\end{proof}

\begin{theorem}[A3.7: local-group nilpotence detection]\label{A:detection}
Let $C$ be a finite category and $k$ any field. For every idempotent
$e\in\operatorname{End}_C(c)$, let $G_e$ be the group of units of
$e\operatorname{End}_C(c)e$. The kernel of the graded ring homomorphism
\[
 H^*(BC;k)\longrightarrow\prod_{(c,e)}H^*(BG_e;k)              \tag{A.3}
\]
is a nilpotent ideal. In particular a positive-degree homogeneous class is
nilpotent if and only if its restriction to every $G_e$ is nilpotent.
Equivalently one may test every subgroup of every local group.
This proves the detection assertion of A3(i) for finite monoids as well.
\end{theorem}
\begin{proof}
\emph{The cohomology ring and the action.}
For clarity, all module diagrams here are covariant. A projective resolution
of the constant diagram has degree-$n$ term
\[
 P_n=\bigoplus_{c_0\xrightarrow{f_1}c_1\longrightarrow\cdots
                 \xrightarrow{f_n}c_n} k\operatorname{Hom}_C(c_n,-).
\]
Its differential deletes end vertices or composes adjacent arrows, using
precomposition on the representable for the last face. At an object $d$
the augmented complex is the simplicial chain complex of $C/d$, whose
terminal object is $1_d$, and is therefore exact. Applying
$\operatorname{Hom}_{kC}(-,\Delta k)$ gives precisely the nerve cochains.
The Alexander--Whitney diagonal on this resolution sends a chain ending
in $h:c_n\to d$ to the sum of its front and back pieces, with the front
piece ending in $hf_n\cdots f_{i+1}:c_i\to d$ and the back piece ending
in $h$. Its associated product is
\[
 (a\smile b)(f_1,\ldots,f_{p+q})
   =a(f_1,\ldots,f_p)b(f_{p+1},\ldots,f_{p+q}).
\]
This identifies the Yoneda product on
$\operatorname{Ext}_{kC}^*(\Delta k,\Delta k)$ with the ordinary cup product.
Indeed the diagonal lifts the identity of the tensor unit, and composing
representatives by this lift is the standard comparison construction of
the Yoneda product. Pointwise tensor over a field is exact, so it descends
to the symmetric monoidal bounded derived category. Tensoring a morphism
of its unit with an object gives the graded central action used below;
acting on the unit gives that same Yoneda, hence cup, product. This also
explains explicitly the ring and restriction identifications used in
Lemma~\ref{A:action-induction}.

\emph{First, the monoid case.}
Write $J$ for the kernel of (A.3). The ring $H^*(BM;k)$ acts by tensoring
on the bounded derived category of finite-dimensional $kM$-modules.
This is a graded central action. Let $\mathcal N$ be the class of objects
annihilated by some power of $J$. This class is closed under shifts,
finite direct sums, direct summands and exact triangles. For the last
assertion, the elementary ghost argument is as follows. In a triangle
$A\to B\to D\to A[1]$, if a homogeneous product $z\in J^s$ kills $B$,
then $z_D$ factors through $D\to A[1]$, by the exact sequence of morphisms
out of the triangle. If $J^r$ kills $A$, composing with any homogeneous
product in $J^r$ gives zero by naturality of the central action. Therefore
$J^{r+s}$ kills $D$. Rotating the triangle gives the other closure cases.
Linearity handles sums of homogeneous products.

By Lemma~\ref{A:action-induction}, $J$ annihilates every
$\operatorname{Ind}_eW$. Descending induction on the finite partially
ordered set of regular $\mathcal J$-classes, using (A.1), puts every simple
module in $\mathcal N$: the kernel in (A.1) has finite length and only
higher-apex composition factors. A composition series then puts every
finite-dimensional module in $\mathcal N$, in particular the trivial
module $k$. Thus $J^Nk=0$ for some integer $N$. Acting on the unit is the
identity identification with $H^*(BM;k)$, so $J^N=0$ as an ideal.

\emph{The finite-category case.}
Here the tensor unit is the constant functor $\Delta k$ and tensor products
are pointwise. To justify the same finite filtration, form the consolidation
monoid $S=\operatorname{Mor}(C)\amalg\{0,1\}$, with a fresh identity and
zero, and multiply noncomposable arrows to zero. A $kC$-module extends to
$S$ on the direct sum of its object values; the fresh zero acts as zero and
$\sum_c1_c$ acts as the identity. Its simple apices, by the proof of
Lemma~\ref{A:apex}, are actual idempotent endomorphisms of $C$: the fresh
zero annihilates it and an actual object identity acts nontrivially, so the
fresh identity cannot be a minimal apex. The induced modules of (A.1)
are again $kC$-modules, because their basis arrows have one target and the
sum of the object identities acts as one.

More explicitly, for an apex $e:c\to c$ their value at $d$ is
$kL_e(d)\otimes_{kG_e}W$, where $L_e(d)$ consists of the arrows
$l:c\to d$ in the $\mathcal L$-class of $e$. The natural transformation
(A.2) at object $d$ is
\[
 l\otimes(w\otimes v)\longmapsto(l\otimes w)\otimes V(l)v,
 \qquad v\in eV(c).
\]
It is well defined, natural and exact by exactly the checks above, now
pointwise, and is an isomorphism on $\Delta k$. Hence the kernel in (A.3)
annihilates all these induced modules. The same apex and composition-series
argument makes a power annihilate $\Delta k$. The bar resolution of
$\Delta k$ identifies its Ext ring, with this action, with $H^*(BC;k)$.
It proves the nilpotent-kernel assertion for $C$.

Finally there are only finitely many idempotents. If all restrictions of
$x$ are nilpotent, one common power $x^a$ belongs to $J$, and then
$x^{aN}=0$. Restriction preserves nilpotence in the other direction.
Testing all subgroups is equivalent, since the full $G_e$ are included.
\end{proof}

\begin{remark}[An explicit uniform exponent]\label{A:detection-exponent}
The proof bounds the power of the entire kernel, not only a power of one
class. For each simple module $S$ choose (A.1), and define recursively
\[
 a(S)=1+\sum_T[K:T]a(T).
\]
The sum runs over the simple composition factors of $K$, with multiplicity;
the strictly larger apices make this a finite recursion. Then $J^{a(S)}$
annihilates $S$, and
$N=\sum_S[\Delta k:S]a(S)$ is a valid nilpotence exponent in (A.3).
There is also a rough bound depending only on the category. Put
$q=|\operatorname{Ob}C|$, $t=|\operatorname{Mor}C|$, and let $r$ be the
number of regular $\mathcal J$-classes of the consolidation monoid
represented by actual arrows. For a nonempty category one may take
\[
 N\leq q\sum_{j=0}^{r-1}t^j.
\]
To check this, a simple $kG_e$-module is a quotient of its regular module,
so has dimension at most $|G_e|$. Freeness of $L_e$ then gives
$\dim\operatorname{Ind}_e(eS)\leq |L_e|\leq t$; hence each kernel has
composition length at most $t$. Induction over at most $r$ successive
apex levels gives the displayed geometric sum. Finally the total vector
space dimension of $\Delta k$ is $q$, so its composition length is at most
$q$. The empty category is vacuous.
\end{remark}

\begin{corollary}[A3.8: finite rational cup length]\label{A:rational-nilpotence}
For every finite category $C$ there is an integer $N$ such that every
product of $N$ positive-degree rational cohomology classes of $BC$ is zero.
In particular neither $K(\mathbb Z,2)$ nor $\Omega S^3$ is the classifying
space of a finite category or finite monoid.
\end{corollary}
\begin{proof}
Averaging over a finite group makes the invariants functor over
$\mathbb Q$ exact, so $H^{>0}(BG_e;\mathbb Q)=0$ for every local group.
The positive-degree ideal therefore lies in the nilpotent kernel of
Theorem~\ref{A:detection}. The rational cohomology of $K(\mathbb Z,2)$ is
$\mathbb Q[u]$, $|u|=2$. The rational cohomology of $\Omega S^3$ is also a
polynomial algebra on a degree-two class: the path-loop fibration, or the
standard tensor-algebra computation $H_*(\Omega\Sigma S^2;\mathbb Q)$,
gives a divided-power cohomology algebra which becomes polynomial over
$\mathbb Q$. Both have arbitrarily long nonzero products, contrary to the
first assertion.
\end{proof}

\begin{corollary}[A3.9: a modular obstruction with infinitely many polynomial generators]\label{A:K22}
The space $K(\mathbb Z/2,2)$ is not homotopy equivalent to the classifying
space of a finite category, and hence not of a finite monoid.
More generally, for any finite category $C$, every polynomial subalgebra
of $H^*(BC;\mathbb F_2)$ has a bounded finite number of algebraically
independent homogeneous generators.
\end{corollary}
\begin{proof}
By Evens's finite-generation theorem, the mod-$2$ cohomology ring of each
finite $G_e$ is a finitely generated algebra. This is Theorem~8.1, page~238,
of \emph{The cohomology ring of a finite group} (1961) \cite{evens1961}, specialized to the
trivial coefficient ring $\mathbb F_2$. Choose homogeneous positive-degree
generators, and let $d$ be the maximum of their numbers over the finitely
many local groups. The cumulative dimensions through degree $n$ of the
product in (A.3) are therefore at most $A(n+1)^d$ for some constant $A$:
monomials of total ordinary exponent at most $n$ span these degrees.

A polynomial subalgebra of $H^*(BC;\mathbb F_2)$ intersects the nilpotent
kernel of (A.3) in zero, since it has no nonzero nilpotent elements. It
therefore injects into that product. If it had $d+1$ algebraically
independent homogeneous generators, and their degrees were at most $b$,
then its cumulative dimensions would be at least
$\binom{\lfloor n/b\rfloor+d+1}{d+1}$, contradicting the preceding
degree-$d$ polynomial bound for large $n$. This proves the general claim.

Serre's computation, \emph{Cohomologie modulo 2 des complexes
d'Eilenberg--MacLane} (1953), Section~2, Theorem~2, page~203 \cite{serre1953}, identifies
$H^*(K(\mathbb Z/2,2);\mathbb F_2)$ with a polynomial algebra. Its
generators include
\[
 u_2,\quad \operatorname{Sq}^{1}u_2,\quad
 \operatorname{Sq}^{2}\operatorname{Sq}^{1}u_2,\quad
 \operatorname{Sq}^{4}\operatorname{Sq}^{2}\operatorname{Sq}^{1}u_2,
 \quad\ldots,
\]
of degrees $2,3,5,9,\ldots$. The nonempty indexing sequences are
admissible of excess one, so the theorem makes these infinitely many
classes algebraically independent. They contradict the general claim.
Both external theorem statements were checked in the original articles.
\end{proof}

\subsection{A5: the bounded-homology assertion without a hypothesis on the fundamental group}

\begin{lemma}[A5.3: a finite proper action associated with the groupoid image]\label{A:groupoid-image}
Let $C$ be a finite connected category and let $\Pi C$ be its groupoid
localization. There is a finite category $D$, with the same objects and
faithful inclusion into $\Pi C$, and a functor $q:C\to D$ inducing an
equivalence of groupoid localizations. All endomorphisms of $D$ are
invertible. If $\Gamma=\pi_1(BC)$, there is a finite proper
$\Gamma$-CW-complex $Y$, of dimension at most the number of isomorphism
classes of objects of $D$ minus one, and a homotopy equivalence
\[
 BD\simeq E\Gamma\times_\Gamma Y.                       \tag{A.4}
\]
Under this equivalence the inclusion of
$H_c=\operatorname{Aut}_D(c)$ corresponds to the map from the classifying
space of a vertex stabilizer of $Y$.
\end{lemma}
\begin{proof}
Define $\operatorname{Hom}_D(c,d)$ to be the image of
$\operatorname{Hom}_C(c,d)$ in $\operatorname{Hom}_{\Pi C}(c,d)$.
Composition is inherited from $\Pi C$, so this is a finite category.
The universal property of localization gives mutually inverse functors
$\Pi C\to\Pi D\to\Pi C$: their composites are the identity on the
generating arrows. A finite submonoid of a group is a subgroup, so every
endomorphism of $D$ is invertible.

Fix a base object $b$. Define $\widetilde D$ to have objects $(c,\gamma)$,
where $\gamma:b\to c$ is an arrow of $\Pi C$, and an arrow
$(c,\gamma)\to(d,\delta)$ whenever $\delta=f\gamma$ for $f\in D(c,d)$.
Such an $f$ is unique because $D\to\Pi C$ is faithful. Thus
$\widetilde D$ is a preorder. The group
$\Gamma=\operatorname{Aut}_{\Pi C}(b)$ acts freely on its objects by
$a(c,\gamma)=(c,\gamma a^{-1})$; its quotient category is $D$.
The nerve projection is a covering: a chain lifts uniquely once a lift
of any one of its vertices is specified. In particular
$B\widetilde D/\Gamma=BD$, and the Borel construction on
$B\widetilde D$ is homotopy equivalent to $BD$.

Let $P$ be the poset obtained from this preorder by identifying isomorphic
objects, and put $Y=BP$. The functor $\widetilde D\to P$ is an
equivalence of categories, is $\Gamma$-equivariant, and induces a homotopy
equivalence on underlying spaces. Its Borel construction therefore also
induces a homotopy equivalence: it is a map of associated fibrations over
$B\Gamma$ with a homotopy equivalence on fibers. This gives (A.4).

A strict chain in $P$ cannot revisit an isomorphism class of objects of
$D$. Indeed a composite between isomorphic objects of the EI category $D$
is invertible, and its inverse lifts to a reverse arrow in $\widetilde D$. Thus
the corresponding endpoints in $\widetilde D$ would already be
isomorphic. This bounds the dimension as stated. Every chain of $P$
lifts to a chain of $\widetilde D$; the latter has only finitely many
$\Gamma$-orbits of chains in each dimension, since these project to
chains in the finite category $D$. Hence $Y$ has finitely many orbits of
cells in total.

The stabilizer of $[(c,\gamma)]\in P$ is the subgroup
$\gamma^{-1}H_c\gamma$ of $\Gamma$. It is finite. A group element
preserving a simplex of the order complex preserves its ordered vertices,
and hence fixes that simplex pointwise. Thus $Y$ is a proper
$\Gamma$-CW-complex in the required sense. Finally the full subcategory
of $\widetilde D$ over the $H_c$-orbit of $(c,\gamma)$ is the chaotic
category on $H_c$: its nerve is contractible, its quotient is $BH_c$,
and it maps to the single vertex $[(c,\gamma)]$ of $P$. This verifies
the asserted compatibility of (A.4) with the local-group maps.
\end{proof}

\begin{theorem}[A5.4: the unrestricted bounded-cohomology obstruction]\label{A:bounded-categories}
Let $C$ be a finite category. If $H^j(BC;\mathbb Z)=0$ for all sufficiently
large $j$, every endomorphism of $C$ is trivial in its fundamental groupoid.
In particular the homomorphism
$G_e\to\pi_1(BC,c)$ is zero for every idempotent
$e\in\operatorname{End}_C(c)$ and every subgroup of its local group.
The same conclusion holds if the integral homology of $BC$ is bounded.
No finiteness, residual finiteness or linearity assumption on $\pi_1(BC)$
is necessary. The quotient construction of Proposition~\ref{A:quotient}
therefore also applies under this bounded-homology hypothesis.
\end{theorem}
\begin{proof}
It suffices to work in one connected component. Use $q:C\to D$,
$\Gamma$, and $Y$ from Lemma~\ref{A:groupoid-image}.
L\"uck--Oliver's vector-bundle existence theorem supplies a finite-rank
complex $\Gamma$-equivariant vector bundle $E\to Y$ for which every
isotropy group acts on the corresponding fiber as a positive multiple
of its regular representation. This is Corollary~2.7 on page~595 of the
published paper \cite{luckoliver2001}; it is Corollary~2.8 in the author's
preprint. Its hypotheses are that $\Gamma$ be discrete, $Y$ be a
finite-dimensional proper $\Gamma$-CW-complex, and the orders of its
isotropy groups be bounded. Each has been checked in
Lemma~\ref{A:groupoid-image}. The original proof, via Theorem~2.6,
chooses regular representations of one common dimension and extends
them after taking a positive number of direct-sum copies; hence the
multiple is indeed positive.

Form the ordinary Borel vector bundle
$E\Gamma\times_\Gamma E\to E\Gamma\times_\Gamma Y$, transport it
along (A.4), and pull it back to $BC$. Denote its rank by $d>0$ and its
rank-zero K-class by $v$. Under the inclusion $BH_c\to BD$ this bundle
is the bundle of a representation $\rho_c$ which is a positive multiple
of the regular $H_c$-representation. This follows from the last
compatibility statement of Lemma~\ref{A:groupoid-image}.

Choose $n$ with $H^j(BC;\mathbb Z)=0$ for $j>n$. The class $v$ has
skeletal filtration at least one, so Lemma~\ref{A:filtration} makes
$v^{n+1}$ vanish on every finite skeleton of $BC$.
Suppose that $H_c\ne1$. Apply Lemma~\ref{A:minimal-group} to the finite
endomorphism monoid $M_c=\operatorname{End}_C(c)$. A maximal group
$G_e=eM_ce$ in its minimal ideal maps onto the group image $H_c$.
There is consequently an element $g\in G_e$ of prime-power order with
nonidentity image in $H_c$: take a lift of a nonidentity element and then
a primary component whose image is nonidentity. Let $T=\langle g\rangle$,
a cyclic $p$-group.

The local-group map $BT\to B\operatorname{Kar}(C)\simeq BC$, followed
by $Bq$, is the group map $BT\to BH_c\to BD$. Indeed $q(e)=1_c$,
and extending $q$ to $\operatorname{Kar}(C)$ sends $(c,e)$ to $c$.
The pullback of $v$ to $BT$ is therefore
$\rho_c|_T-d\in R(T)$ under the completion map. Since the pullback
of $v^{n+1}$ vanishes on all finite skeleta of $BT$, Atiyah's Theorem
(7.2) and Proposition (6.11) \cite{atiyah1961} imply
\[
 (\rho_c|_T-d)^{n+1}=0\qquad\hbox{in }R(T).
\]
But the character of $\rho_c$ at the nonidentity image of $g$ is zero.
Evaluation at $g$ makes the left side $(-d)^{n+1}\ne0$, a contradiction.
Thus every $H_c$ is trivial, which says exactly that every endomorphism
of $C$ dies in $\Pi C$.

Finally, the nerve of a finite category has finitely many simplices in
each degree, so its integral homology is finitely generated degreewise.
The universal coefficient exact sequence converts bounded integral
homology into bounded integral cohomology. This proves the last assertion.
\end{proof}

\begin{corollary}[A5.5: the prime-local version for arbitrary finite categories]
\label{A:bounded-categories-p}
Let $C$ be a finite category and $p$ a prime. If
$H_n(BC;\mathbb F_p)=0$ for all sufficiently large $n$, then the image of
each endomorphism monoid in the fundamental groupoid has order prime to
$p$. Equivalently, every local finite $p$-subgroup has trivial image in
the fundamental groupoid. In particular, if any such image has order
divisible by $p$, then $H_n(BC;\mathbb F_p)\ne0$ for arbitrarily large $n$.
No hypothesis on the ambient fundamental group is needed.
\end{corollary}
\begin{proof}
Every degree of the nerve has finitely many simplices, so the integral
homology and cohomology groups are finitely generated. The universal
coefficient sequences therefore turn bounded mod-$p$ homology into bounded
cohomology with coefficients in $\mathbb Z_{(p)}$.

Use the same finite groupoid-image category $D$, finite proper
$\Gamma$-CW complex $Y$, and regular-isotropy vector bundle as in
Theorem~\ref{A:bounded-categories}. Its Borel bundle, pulled back to $BC$,
has positive rank $d$ and reduced class $v$. The localized form of
Lemma~\ref{A:filtration} says that some power of $v$ restricts to zero
on every finite skeleton in $K^0(-)\otimes\mathbb Z_{(p)}$.

Suppose that one of the finite group images $H_c$ has order divisible
by $p$. Its surjection from the minimal-ideal group $G_e$ of
$\operatorname{End}_C(c)$ supplies a cyclic $p$-subgroup $T\le G_e$
whose image in $H_c$ is nontrivial. Restriction of the bundle to $BT$
is the restriction of a positive multiple of the regular representation
of $H_c$. As in the localized argument in
Theorem~\ref{A:representation}, reduced $K^0$ of the finite odd-dimensional
lens-space skeleta of $BT$ is a finite $p$-group. Localization is injective
on these groups, so the same power of $v|_{BT}$ vanishes integrally on
all these skeleta. Atiyah completion and its injectivity for cyclic
$p$-groups imply
\[
   ([\rho|_T]-d)^N=0\quad\hbox{in }R(T)
\]
for some $N>0$. On an element with nontrivial image in $H_c$, the left
side has character value $(-d)^N$, a contradiction. Thus every $H_c$
has order prime to $p$.

Finally a finite group image has order divisible by $p$ if and only if
the image of some local finite $p$-subgroup is nontrivial: the forward
implication follows by lifting a Sylow subgroup through the surjection
from $G_e$, and the reverse implication is immediate. This proves the
equivalent formulation and its contrapositive.
\end{proof}
